% 07-status-error.tex - Status and Error Handling

\chapter{Status and Error Handling}
\label{sec:status-error}

This section describes \SC{}'s system states, state transitions, and error reporting mechanisms.

\section{System States}

\SC{} reports its operational status via the \mibentry{SUMMARY} MIB entry as a 7-character string (Table~\ref{tab:system-states}).

\begin{table}[htbp]
    \centering
    \caption{System Status Values}
    \label{tab:system-states}
    \begin{tabular}{lp{9cm}}
        \toprule
        \textbf{Status} & \textbf{Description} \\
        \midrule
        \status{BOOTING} & System is initializing (INI in progress) \\
        \status{NORMAL } & System operating normally \\
        \status{SHUTDWN} & System has been shut down \\
        \bottomrule
    \end{tabular}
\end{table}

\section{State Transitions}

The state transition diagram is shown in Figure~\ref{fig:state-transitions}.

\begin{figure}[h]
    \centering
    \begin{verbatim}
                    +----------+
                    | BOOTING  |<----+
                    +----------+     |
                         |           |
                    INI complete     |
                         |           |
                         v           |
                    +----------+     |
                    | NORMAL   |     |
                    +----------+     |
                         |           |
                    SHT command      |
                         |           |
                         v           |
                    +----------+     |
                    | SHUTDWN  |-----+
                    +----------+  INI command
    \end{verbatim}
    \caption{State Transition Diagram}
    \label{fig:state-transitions}
\end{figure}

\subsection{State Transition Events}

\begin{itemize}
    \item \textbf{BOOTING $\rightarrow$ NORMAL}: INI sequence completes successfully; all threads started
    \item \textbf{NORMAL $\rightarrow$ SHUTDWN}: SHT command received; state transitions immediately, then all threads stopped
    \item \textbf{SHUTDWN $\rightarrow$ BOOTING}: INI command received; initialization begins
\end{itemize}

\section{Command Exit Codes}

Commands return the exit codes listed in Table~\ref{tab:exit-codes} to indicate success or failure reason.

\begin{table}[htbp]
    \centering
    \caption{Command Exit Codes}
    \label{tab:exit-codes}
    \begin{tabular}{cp{9cm}}
        \toprule
        \textbf{Code} & \textbf{Description} \\
        \midrule
        \hex{00} & Process accepted without error \\
        \hex{01} & Invalid arguments to command (e.g., unknown \DR{} name or unknown copy ID) \\
        \hex{02} & Other error running command (e.g., error queuing copy or delete) \\
        \hex{03} & Blocking operation in progress (INI or SHT is active) \\
        \hex{04} & Subsystem needs to be initialized (system is in SHUTDWN state) \\
        \bottomrule
    \end{tabular}
\end{table}

\section{Error Response Format}

When a command is rejected, the response data contains the exit code and explanation:

\begin{lstlisting}[caption={Error Response Format}]
0x<CODE>! <LASTLOG MESSAGE>
\end{lstlisting}

Example responses:
\begin{itemize}
    \item \code{0x01! SCP: Invalid arguments to command - unknown data recorder DR7}
    \item \code{0x03! INI: Blocking operation in progress - SHT is active and blocking}
    \item \code{0x04! SCP: Subsystem needs to be initialized}
\end{itemize}

\section{Process Locking}

\SC{} maintains a process list to prevent conflicting operations:

\begin{itemize}
    \item \textbf{INI}: Blocks other INI and SHT commands
    \item \textbf{SHT}: Blocks other INI and SHT commands
\end{itemize}

The \code{activeProcess} list tracks currently running blocking operations. Data management commands (\cmd{SCP}, \cmd{SRM}, \cmd{PAU}, \cmd{RES}, \cmd{SCN}) do not block each other but require the system to be in the \status{NORMAL} state.

\section{Email Notifications}

\SC{} sends email notifications for failed transfers. The notification system:

\begin{itemize}
    \item Runs daily alongside the purge operations
    \item Collects all failed items for each \DR{} queue
    \item Generates a report listing each failed file with its path, size, and failure reason
    \item Sends the report via SMTP with TLS using the sender address and credentials from the \code{email} section of \file{defaults.json} (resolved via \code{lwa\_auth})
    \item The recipient list is currently hardcoded in the source
    \item After notification, failed records are purged from the database
\end{itemize}
